% =============================================================================
\section{Full Proofs}
\label{app:proofs}
% =============================================================================

This appendix provides complete proofs for Theorems~1--4 and Propositions~1--4
stated in the main text.

\subsection{Proof of Proposition~\ref{prop:monotonicity}: Monotonicity and Convexity}

\begin{proof}
From the cost function $C(s, n) = nB + S(2sn + n^2)/2$:

\textbf{(i) Monotonicity in $s$.}
\begin{equation}
  \frac{\partial C}{\partial s} = \frac{\partial}{\partial s}\left[nB + \frac{S(2sn + n^2)}{2}\right] = \frac{S \cdot 2n}{2} = Sn
\end{equation}
Since $S > 0$ and $n > 0$ by assumption, $\partial C / \partial s = Sn > 0$.

\textbf{(ii) Monotonicity in $n$.}
\begin{equation}
  \frac{\partial C}{\partial n} = B + \frac{S(2s + 2n)}{2} = B + S(s + n)
\end{equation}
Since $B > 0$, $S > 0$, and $s, n \geq 0$ with $n > 0$ for a nontrivial
purchase, $\partial C / \partial n = B + S(s + n) > 0$.

Note that $\partial C / \partial n = P(s + n)$, the price at the upper
boundary of the purchase interval. This confirms that the marginal cost of the
last infinitesimal unit equals the bonding curve price at the post-purchase
supply level.

\textbf{(iii) Convexity in $n$.}
\begin{equation}
  \frac{\partial^2 C}{\partial n^2} = \frac{\partial}{\partial n}[B + S(s + n)] = S
\end{equation}
Since $S > 0$, the cost function is strictly convex in $n$. This implies that
the marginal cost is strictly increasing: each additional token costs more than
the previous one. \qed
\end{proof}

\subsection{Proof of Theorem~\ref{thm:curve-comparison}: Quadratic vs.\ Linear vs.\ Exponential}

\begin{proof}
We prove each part in sequence.

\textbf{Setup.} Calibrate the three curves to produce identical total cost at
reference supply $q^*$. Let $C_{\text{lin}}(q^*) = C_{\text{quad}}(q^*) =
C_{\text{exp}}(q^*)$. This yields:
\begin{align}
  C_{\text{lin}}(q^*) &= B_1 \cdot q^* \\
  C_{\text{quad}}(q^*) &= B_2 \cdot q^* + \frac{S_2 (q^*)^2}{2} \\
  C_{\text{exp}}(q^*) &= \frac{B_3}{\alpha}(e^{\alpha q^*} - 1)
\end{align}

Setting these equal determines the calibration relationships.

\textbf{(i) Price ordering for $q > q^*$.}

The exponential price function is $P_{\text{exp}}(q) = B_3 e^{\alpha q}$,
which grows without bound. The quadratic price function is $P_{\text{quad}}(q)
= B_2 + S_2 q$, which is linear. The constant price function is
$P_{\text{lin}}(q) = B_1$.

For sufficiently large $q$, the exponential dominates the linear, which
dominates the constant:
\begin{equation}
  \lim_{q \to \infty} \frac{P_{\text{exp}}(q)}{P_{\text{quad}}(q)} = \lim_{q \to \infty} \frac{B_3 e^{\alpha q}}{B_2 + S_2 q} = +\infty
\end{equation}

By the calibration, at $q = q^*$ the ordering may not hold, but for $q > q^*$
with $q - q^*$ sufficiently large, the exponential exceeds the quadratic which
exceeds the constant. More precisely, since $P_{\text{exp}}''(q) = \alpha^2 B_3
e^{\alpha q} > 0 = P_{\text{quad}}''(q)$ and both are continuous with
calibrated crossing at $q^*$, the exponential lies above the quadratic for
$q > q^*$.

\textbf{(ii) Surplus ratio.}

The early-buyer surplus for the quadratic curve is:
\begin{equation}
  \Delta_{\text{quad}}^{\text{early}} = P_{\text{quad}}(q^*) - P_{\text{quad}}(0) = S_2 q^*
\end{equation}

The late-buyer cost premium for $q > q^*$ is:
\begin{equation}
  \Delta_{\text{quad}}^{\text{late}} = P_{\text{quad}}(q) - P_{\text{quad}}(q^*) = S_2 (q - q^*)
\end{equation}

The ratio is:
\begin{equation}
  R_{\text{quad}} = \frac{S_2 q^*}{S_2(q - q^*)} = \frac{q^*}{q - q^*}
\end{equation}

For the exponential curve:
\begin{align}
  \Delta_{\text{exp}}^{\text{early}} &= B_3(e^{\alpha q^*} - 1) \\
  \Delta_{\text{exp}}^{\text{late}} &= B_3(e^{\alpha q} - e^{\alpha q^*})
\end{align}

The ratio is:
\begin{equation}
  R_{\text{exp}} = \frac{e^{\alpha q^*} - 1}{e^{\alpha q} - e^{\alpha q^*}} = \frac{e^{\alpha q^*} - 1}{e^{\alpha q^*}(e^{\alpha(q - q^*)} - 1)}
\end{equation}

For large $q - q^*$:
\begin{equation}
  R_{\text{exp}} \approx \frac{e^{\alpha q^*} - 1}{e^{\alpha q^*} \cdot e^{\alpha(q - q^*)}} \to 0
\end{equation}

while $R_{\text{quad}} = q^* / (q - q^*)$ decays as $O(1/q)$. Since the
exponential ratio decays as $O(e^{-\alpha q})$, which is faster than
$O(1/q)$, we have $R_{\text{quad}} > R_{\text{exp}}$ for sufficiently large
$q$, confirming inequality~\eqref{eq:surplus-ratio}.

\textbf{(iii) Goldilocks property.}

Accessibility requires $C(q) < \infty$ for all finite $q$. This holds for the
constant ($C = B_1 q$), quadratic ($C = B_2 q + S_2 q^2 / 2$), and
exponential ($C = (B_3 / \alpha)(e^{\alpha q} - 1)$) curves. However, the
exponential cost grows as $e^{\alpha q}$, making participation at large $q$
economically prohibitive even though mathematically finite. We formalize this
by defining \emph{practical accessibility} as $C(q) \leq \kappa \cdot C(q^*)$
for a constant $\kappa$ (e.g., $\kappa = 10$). The quadratic curve satisfies
this for $q \leq q^* \sqrt{2\kappa}$, while the exponential curve satisfies
this only for $q \leq q^* + \ln(\kappa) / \alpha$, which is much more
restrictive for typical $\alpha$.

Scarcity requires $P'(q) > 0$. The constant curve fails ($P' = 0$). The
quadratic satisfies $P'(q) = S > 0$. The exponential satisfies $P'(q) =
\alpha B_3 e^{\alpha q} > 0$ but with the excessive growth noted above.

The quadratic curve is the unique member of the polynomial family $P(q) = B +
Sq^k$ with $k = 1$ that satisfies both $P'(q) > 0$ (scarcity) and
$P''(q) = 0$ (linear price growth, ensuring practical accessibility). \qed
\end{proof}

\subsection{Proof of Theorem~\ref{thm:folk}: Cooperation via Lock Commitment}

\begin{proof}
We apply the folk theorem for finitely repeated games with known
horizon~\citep{fudenberg1991game}.

Consider the grim-trigger strategy: ``Stake in every period; if the opponent
ever Sells, Sell in all subsequent periods.'' Under this strategy profile, the
payoff from perpetual cooperation is:
\begin{equation}
  V_{\text{coop}} = \sum_{t=0}^{T-1} \delta^t (r + g - c) = (r + g - c) \cdot \frac{1 - \delta^T}{1 - \delta}
\end{equation}

The best deviation payoff (Sell in period 0, then mutual Sell thereafter) is:
\begin{equation}
  V_{\text{dev}} = \ell + \sum_{t=1}^{T-1} \delta^t (\ell - d) = \ell + (\ell - d) \cdot \delta \cdot \frac{1 - \delta^{T-1}}{1 - \delta}
\end{equation}

For cooperation to be sustained, we need $V_{\text{coop}} \geq V_{\text{dev}}$:
\begin{align}
  (r + g - c) \cdot \frac{1 - \delta^T}{1 - \delta} &\geq \ell + (\ell - d) \cdot \delta \cdot \frac{1 - \delta^{T-1}}{1 - \delta}
\end{align}

For $T \to \infty$ (or equivalently, for $T$ large enough that $\delta^T
\approx 0$), this simplifies to:
\begin{equation}
  \frac{r + g - c}{1 - \delta} \geq \ell + \frac{(\ell - d)\delta}{1 - \delta}
\end{equation}

Multiplying through by $(1 - \delta)$:
\begin{equation}
  r + g - c \geq \ell(1 - \delta) + (\ell - d)\delta = \ell - d\delta
\end{equation}

Rearranging:
\begin{equation}
  \delta \geq \frac{\ell - r - g + c}{\ell - r + c - (d - g) + g} = \frac{\ell - r - g + c}{d}
\end{equation}

Wait---let us redo this more carefully. From the original condition:
\begin{equation}
  r + g - c \geq \ell - \ell\delta + \ell\delta - d\delta = \ell - d\delta
\end{equation}

So: $r + g - c - \ell + d\delta \geq 0$, giving:
\begin{equation}
  \delta \geq \frac{\ell - (r + g - c)}{d}
\end{equation}

However, the \veCORTEX{} lock mechanism provides a stronger result: during
the lock period, the Sell action is \emph{mechanically unavailable}. The
strategy space is reduced to $\{\text{Stake}\}$ for the duration of the lock.
This is equivalent to setting the deviation payoff to $-\infty$ during the
lock period, making cooperation the unique available strategy regardless of
the discount factor. The folk theorem threshold~\eqref{eq:folk-threshold}
then serves as the condition for voluntary re-locking after the initial lock
expires. \qed
\end{proof}

\subsection{Proof of Proposition~\ref{prop:creator-wealth}: Vesting Maximizes Expected Wealth}

\begin{proof}
The creator's expected wealth at day $d$ is:
\begin{equation}
  W(d) = V(d) \cdot P(d) + \sum_{t=d}^{\infty} \beta^{t-d} \cdot f_c \cdot TV(t)
\end{equation}

Under the vesting schedule $V(d) = 0.01 A_c + 0.99 A_c \min(d, 180)/180$,
the creator cannot sell more than $V(d)$ at day $d$. Consider a creator
choosing effort level $e$ at cost $\psi(e)$ (convex, increasing). Both $P$
and $TV$ depend on $e$:
\begin{equation}
  P(d; e) = P_0 + \phi_P(e), \quad TV(t; e) = TV_0 + \phi_{TV}(e)
\end{equation}
where $\phi_P, \phi_{TV}$ are increasing in $e$ (effort improves agent quality,
which increases token price and trading volume).

The creator maximizes:
\begin{equation}
  \max_e \; W(d; e) - \psi(e)
\end{equation}

The first-order condition is:
\begin{equation}
  V(d) \cdot \phi_P'(e) + \sum_{t=d}^{\infty} \beta^{t-d} f_c \phi_{TV}'(e) = \psi'(e)
\end{equation}

Now compare two regimes:

\textbf{Regime 1 (Vesting):} At $d = 0$, $V(0) = 0.01 A_c$. The creator
cannot liquidate and must rely primarily on the fee stream. The optimal effort
$e^*_V$ satisfies:
\begin{equation}
  0.01 A_c \cdot \phi_P'(e^*_V) + \frac{f_c \phi_{TV}'(e^*_V)}{1 - \beta} = \psi'(e^*_V)
\end{equation}

\textbf{Regime 2 (No vesting):} At $d = 0$, $V(0) = A_c$. The creator can
immediately sell $A_c \cdot P_0$ and exit. The optimal effort $e^*_N$ satisfies:
\begin{equation}
  A_c \cdot \phi_P'(e^*_N) + \frac{f_c \phi_{TV}'(e^*_N)}{1 - \beta} = \psi'(e^*_N)
\end{equation}

In Regime 2, the creator has a profitable deviation: set $e = 0$, sell $A_c$
at $P_0$, and exit. This yields $W = A_c \cdot P_0$, which may exceed the
effort-adjusted long-run value if $\psi(e^*_N)$ is large.

In Regime 1, the deviation to $e = 0$ yields only $W = 0.01 A_c \cdot P_0$,
which is 99\% lower. The creator is thus incentivized to choose $e > 0$ to
increase $P(d)$ for $d > 0$ (when more tokens vest) and to build the fee
stream. The vesting constraint effectively forces the creator to internalize
the future, converting a short-horizon liquidation problem into a long-horizon
optimization problem. \qed
\end{proof}

\subsection{Proof of Theorem~\ref{thm:separating}: Separating Equilibrium}

\begin{proof}
Following the Spence signaling framework~\citep{spence1973job}, define:
\begin{itemize}
  \item Types: $\theta \in \{\theta_L, \theta_H\}$ with $\theta_H > \theta_L$.
  \item Signal: Platform choice $\sigma \in \{P, A\}$ (PURECORTEX with vesting,
    or Platform A without vesting).
  \item Cost of signal: $c(\theta, \sigma)$ where $c(\theta, P) > c(\theta, A)$
    for all $\theta$ (vesting is costlier than no vesting).
\end{itemize}

The single-crossing condition requires:
\begin{equation}
  \frac{\partial c}{\partial \sigma}\bigg|_{\theta_H} < \frac{\partial c}{\partial \sigma}\bigg|_{\theta_L}
\end{equation}

Interpreted: the marginal cost of the vesting signal (choosing PURECORTEX over
Platform A) is lower for high-quality creators than for low-quality creators.

This holds because:
\begin{itemize}
  \item A high-quality creator ($\theta_H$) who continues developing their
    agent expects $P(d)$ to increase or remain stable during the 180-day
    vesting. The effective cost of vesting is $c_H = A_c \cdot [P(0) - P(180; \theta_H)]$, which is negative (a benefit) if the token appreciates.
  \item A low-quality creator ($\theta_L$) who plans to abandon the agent
    expects $P(d)$ to decline. The effective cost is $c_L = A_c \cdot [P(0) - P(180; \theta_L)]$, which is positive and potentially large.
\end{itemize}

For a separating equilibrium:
\begin{align}
  \theta_H: &\quad \theta_H \cdot \Delta P - c_H \geq 0 \quad \text{(prefer PURECORTEX)} \\
  \theta_L: &\quad 0 \geq \theta_L \cdot \Delta P - c_L \quad \text{(prefer Platform A)}
\end{align}

where $\Delta P$ is the price premium buyers assign to PURECORTEX-launched
agents. This gives the separating condition:
\begin{equation}
  \frac{c_L}{\theta_L} \geq \Delta P \geq \frac{c_H}{\theta_H}
\end{equation}

Since $c_H < c_L$ (vesting is less costly for high-quality creators) and
$\theta_H > \theta_L$, we have $c_H / \theta_H < c_L / \theta_L$, so the
interval is nonempty and a separating equilibrium exists.

In equilibrium, buyers correctly infer $\theta_H$ from the PURECORTEX signal
and $\theta_L$ from the Platform A signal, pricing agent tokens accordingly.
The vesting mechanism serves as a credible commitment device that only
high-quality creators are willing to undertake. \qed
\end{proof}

\subsection{Proof of Theorem~\ref{thm:pareto}: Pareto Improvement via Governance}

\begin{proof}
See the proof in Section~\ref{subsec:senator-mechanism}. We provide additional
detail here.

Let $o_0$ be the current parameter configuration and $o^*$ be the
Senator-proposed configuration with $W(o^*) > W(o_0) + \epsilon$.

The welfare function $W = \alpha \cdot \text{TVL} + \beta \cdot N + \gamma
\cdot \text{ret} - \delta \cdot \text{risk}$ is a weighted sum of ecosystem
health metrics. We need to show that an increase in $W$ implies a Pareto
improvement.

Decompose the agent population into stakers ($\mathcal{S}$), creators
($\mathcal{C}$), and traders ($\mathcal{T}$). Each group's utility depends
on different welfare components:
\begin{align}
  u_i^{\mathcal{S}} &= h_1(\text{TVL}, \text{ret}) \quad \text{(stakers benefit from TVL and retention)} \\
  u_i^{\mathcal{C}} &= h_2(N, \text{TVL}) \quad \text{(creators benefit from agent count and TVL)} \\
  u_i^{\mathcal{T}} &= h_3(\text{TVL}, -\text{risk}) \quad \text{(traders benefit from TVL and low risk)}
\end{align}

where $h_1, h_2, h_3$ are increasing in each argument. If $W(o^*) > W(o_0)$,
then at least one welfare component has increased. By the monotonicity of
group utilities in welfare components, the corresponding group is strictly
better off. By the constraint that the Senator only proposes changes that do
not decrease any individual welfare component (enforced by the constitutional
check), no group is worse off.

More formally, the Senator's proposal algorithm operates as follows:
\begin{enumerate}
  \item Compute $W(o)$ for each candidate configuration $o$.
  \item Filter: retain only configurations where each welfare component is
    weakly greater than under $o_0$.
  \item Select $o^* = \argmax_o W(o)$ from the filtered set.
\end{enumerate}

Step 2 ensures the Pareto condition: no welfare component decreases, so no
group is worse off. Step 3 ensures welfare maximization among Pareto-improving
alternatives. \qed
\end{proof}

\subsection{Proof of Proposition~\ref{prop:delegation}: Rationality of Delegation}

\begin{proof}
See the proof in Section~\ref{subsec:delegation-game}. The result follows
from comparing expected utilities under direct voting and delegation, with the
key insight that delegation to a better-informed delegate strictly dominates
direct voting whenever the delegate's informational advantage exceeds the
agency cost (which is bounded by the delegate's own incentive alignment via
\veCORTEX{} staking). \qed
\end{proof}
